% Especificación del Analizador Léxico (Lenguaje Reducido)
\documentclass[12pt]{article}

\usepackage[T1]{fontenc}
\usepackage[utf8]{inputenc}
\usepackage[spanish]{babel}
\usepackage{geometry}
\usepackage{hyperref}
\usepackage{enumitem}
\usepackage{verbatim}

\geometry{margin=2.5cm}

\title{Especificación del Analizador Léxico}
\author{}
\date{}

\begin{document}
\maketitle

\section{Introducción}
Este proyecto implementa la primera fase de un compilador: el \textbf{análisis léxico}. El objetivo es transformar el código fuente del lenguaje reducido en una secuencia de \textbf{tokens} que pueda consumir la fase sintáctica (Bison). Un compilador depende de esta etapa para:
\begin{itemize}[nosep]
  \item normalizar la entrada (por ejemplo, ignorar espacios y comentarios),
  \item reconocer patrones válidos (identificadores, números, operadores, etc.),
  \item reportar errores tempranos cuando se detectan caracteres o lexemas inválidos.
\end{itemize}

\subsection*{Alcance del proyecto}
\textbf{Sí hace:}
\begin{itemize}[nosep]
  \item Reconoce tokens para: funciones, declaraciones (con posible mutabilidad), asignaciones, expresiones aritmético-lógicas, condicionales \texttt{if}/\texttt{else} y \texttt{return}.
  \item Reconoce: identificadores, palabras reservadas, literales numéricos (enteros y reales), operadores, delimitadores, espacios en blanco y comentarios.
  \item Maneja errores léxicos: caracteres no reconocidos y lexemas inválidos.
  \item Produce tokens consistentes para el parser (Bison).
\end{itemize}

\textbf{No hace:}
\begin{itemize}[nosep]
  \item No soporta ciclos, patrón \texttt{match}, arreglos, cadenas, ni literales de carácter.
  \item No realiza verificación semántica (tipos, ámbitos, etc.).
\end{itemize}

\section{Diseño del Analizador Léxico}
La implementación objetivo se plantea con \textbf{Flex} como generador de analizador léxico: a partir de expresiones regulares, Flex construye un autómata determinista (DFA) equivalente. El lexer se integra con \textbf{Bison} retornando códigos de token y, cuando corresponda, adjuntando el lexema como valor semántico.

\subsection{Alfabeto de entrada}
Se asume entrada en texto plano (ASCII extendido o UTF-8, pero el lenguaje solo define tokens sobre el subconjunto ASCII):
\begin{itemize}[nosep]
  \item Letras: \texttt{A--Z}, \texttt{a--z}
  \item Dígitos: \texttt{0--9}
  \item Guion bajo: \texttt{\_}
  \item Operadores: \texttt{+ - * / \% = ! < > \& |}
  \item Delimitadores: \texttt{( ) \{ \} ; , :}
  \item Punto: \texttt{.} (solo válido dentro de reales)
  \item Espacios en blanco: espacio y tabulación
  \item Saltos de línea: \texttt{\textbackslash n} y \texttt{\textbackslash r\textbackslash n}
\end{itemize}

\subsection{Conjunto de tokens a reconocer}
El lexer debe reconocer (como mínimo) el siguiente conjunto:
\begin{itemize}[nosep]
  \item \textbf{Identificadores}: nombres de variables/funciones.
  \item \textbf{Palabras reservadas}: \texttt{fn}, \texttt{let}, \texttt{mut}, \texttt{if}, \texttt{else}, \texttt{return}, \texttt{true}, \texttt{false}, \texttt{i32}, \texttt{f64}, \texttt{bool}.
  \item \textbf{Números}: enteros decimales y reales decimales con exponente opcional.
  \item \textbf{Operadores}: aritméticos \texttt{+ - * / \%}, comparación \texttt{== != < > <= >=}, lógicos \texttt{\&\& || !}, asignación \texttt{=}.
  \item \textbf{Delimitadores}: \texttt{( ) \{ \} ; , :}.
  \item \textbf{Espacios/Comentarios}: se consumen y no se entregan al parser (salvo que el proyecto exija tokens explícitos de espacio; por defecto se ignoran).
\end{itemize}

\subsection{Definiciones y expresiones regulares}
A continuación se listan las expresiones regulares normativas para Flex.

\subsubsection*{Identificadores}
\begin{verbatim}
ID  [_a-zA-Z](?:[_a-zA-Z0-9]|_(?!_))*
\end{verbatim}
Regla adicional: un identificador \textbf{no} puede contener \texttt{__} consecutivos.

\subsubsection*{Números}
Importante: el signo \texttt{+}/\texttt{-} se tokeniza como operador; la negación/unario se resuelve en la fase sintáctica.

Entero:
\begin{verbatim}
INT  [0-9]+
\end{verbatim}

Real (con exponente opcional):
\begin{verbatim}
REAL  [0-9]+\.[0-9]+([eE][+-]?[0-9]+)?
\end{verbatim}

Nota: un \texttt{.} aislado (o un \texttt{.} no seguido de dígitos) se considera \textbf{error léxico}.

\subsubsection*{Comentarios}
\begin{verbatim}
LINE_COMMENT   //[^\r\n]*
BLOCK_COMMENT  /\*([^*]|\*+[^*/])*\*+/
\end{verbatim}
Los comentarios se consumen (se ignoran) y deben actualizar correctamente línea/columna.

\subsubsection*{Espacios en blanco}
\begin{verbatim}
WS  [ \t]+
NL  (\r\n|\n)
\end{verbatim}
\texttt{WS} y \texttt{NL} se consumen; \texttt{NL} incrementa el contador de línea y reinicia columna.

\subsubsection*{Operadores y delimitadores}
Se recomienda priorizar primero los tokens multi-caracter y luego los de un carácter.
\begin{verbatim}
==  !=  <=  >=  &&  ||
+  -  *  /  %  =  !  <  >
;  ,  :  (  )  {  }
\end{verbatim}

\subsection{Manejo de errores léxicos}
El lexer debe:
\begin{itemize}[nosep]
  \item Reportar \textbf{carácter no reconocido} (cualquier símbolo fuera del alfabeto definido) indicando línea y columna.
  \item Reportar \textbf{token inválido} cuando un patrón casi-válido viola una regla (por ejemplo, identificador con \texttt{__}, o punto aislado).
  \item Continuar el análisis en lo posible (estrategia de recuperación mínima: consumir el carácter/lexema inválido y seguir).
\end{itemize}

\section{Diseño del Analizador Sintáctico y Semántico}
La fase sintáctica se implementa con \textbf{Bison} (Yacc). El analizador sintáctico consume la secuencia de tokens producida por el lexer (Flex) y construye una representación estructurada del programa (por ejemplo, AST) o ejecuta acciones semánticas durante el parseo.

\subsection{Decisiones de diseño (sintáctico)}
\begin{itemize}[nosep]
  \item \textbf{Gramática no ambigua}: se definen precedencias y asociatividades para eliminar ambigüedades de expresiones (aritméticas, relacionales y lógicas).
  \item \textbf{Integración con el lexer}: el lexer retorna tokens como \texttt{TOKEN\_KW\_FN}, \texttt{TOKEN\_IDENTIFIER}, \texttt{TOKEN\_NUMBER}, etc. Bison debe definir estos tokens (en el header generado) o mapearlos de forma consistente.
  \item \textbf{Recuperación básica}: se usa el token especial \texttt{error} de Bison para sincronizar el análisis al encontrar errores (por ejemplo, hasta el siguiente \texttt{;} o \texttt{\}}).
\end{itemize}

\subsection{Definición de gramática (BNF/EBNF)}
La siguiente gramática es una guía (EBNF) alineada con el lenguaje reducido descrito en \texttt{tokens.md} y \texttt{reglas.md}. El objetivo es soportar: funciones, declaraciones \texttt{let} (con \texttt{mut} opcional), asignaciones, sentencias \texttt{return}, condicionales \texttt{if/else}, y expresiones.

\begin{verbatim}
programa        -> ( funcion )* EOF

funcion         -> 'fn' IDENT '(' params_opt ')' bloque
params_opt      -> params | epsilon
params          -> param ( ',' param )*
param           -> IDENT ':' tipo

bloque          -> '{' ( sentencia )* '}'

sentencia       -> let_stmt ';'
               | return_stmt ';'
               | if_stmt
               | expr_stmt ';'

let_stmt        -> 'let' mut_opt IDENT tipo_opt init_opt
mut_opt         -> 'mut' | epsilon
tipo_opt        -> ':' tipo | epsilon
init_opt        -> '=' expr | epsilon

return_stmt     -> 'return' expr_opt
expr_opt        -> expr | epsilon

if_stmt         -> 'if' expr bloque else_opt
else_opt        -> 'else' bloque | epsilon

expr_stmt       -> expr

expr            -> asignacion
asignacion      -> IDENT '=' asignacion | log_or

log_or          -> log_and ( '||' log_and )*
log_and         -> igualdad ( '&&' igualdad )*
igualdad        -> rel ( ( '==' | '!=' ) rel )*
rel             -> suma ( ( '<' | '<=' | '>' | '>=' ) suma )*
suma            -> mult ( ( '+' | '-' ) mult )*
mult            -> unario ( ( '*' | '/' | '%' ) unario )*
unario          -> ( '!' | '+' | '-' ) unario | primario

primario        -> NUMBER
               | 'true'
               | 'false'
               | IDENT
               | IDENT '(' args_opt ')'
               | '(' expr ')'

args_opt        -> args | epsilon
args            -> expr ( ',' expr )*

tipo            -> 'i32' | 'f64' | 'bool'
\end{verbatim}

\subsection{Acciones semánticas: revisión de tipos (mínimo)}
La revisión de tipos se realiza como verificación semántica básica durante el parseo o en un pase posterior sobre el AST. Reglas mínimas recomendadas:
\begin{itemize}[nosep]
  \item \textbf{Declaraciones}: \texttt{let x: T = expr;} requiere que el tipo inferido/obtenido de \texttt{expr} sea compatible con \texttt{T}.
  \item \textbf{Asignaciones}: \texttt{x = expr;} requiere que \texttt{x} exista en la tabla de símbolos y que el tipo de \texttt{expr} sea compatible.
  \item \textbf{Operadores}:
    \begin{itemize}[nosep]
      \item Aritméticos (\texttt{+ - * / %}): operandos numéricos; resultado numérico.
      \item Relacionales (\texttt{< <= > >=}): operandos numéricos; resultado \texttt{bool}.
      \item Igualdad (\texttt{== !=}): operandos de tipos compatibles; resultado \texttt{bool}.
      \item Lógicos (\texttt{\&\& || !}): operandos booleanos; resultado \texttt{bool}.
    \end{itemize}
  \item \textbf{return}: dentro de una función, el valor retornado debe ser una expresión válida (y si se modela tipo de retorno, debe ser compatible).
  \item \textbf{if/else}: la condición debe ser \texttt{bool}.
\end{itemize}

Una implementación típica define un enum de tipos (\texttt{TYPE\_I32}, \texttt{TYPE\_F64}, \texttt{TYPE\_BOOL}, \texttt{TYPE\_ERROR}) y propaga tipos en las reglas de expresión mediante atributos semánticos. Cuando se detecta incompatibilidad, se reporta un error semántico con línea/columna y se propaga \texttt{TYPE\_ERROR} para evitar cascadas.

\end{document}
